% Generated draft: requires mathematical and source review.
\documentclass[11pt]{article}
\usepackage[margin=25mm]{geometry}
\usepackage{fontspec}
\setmainfont{Latin Modern Roman}
\usepackage{amsmath,amssymb,mathtools,unicode-math}
\setmathfont{Latin Modern Math}
\usepackage{enumitem,needspace,xurl,hyperref}
\setlength{\parindent}{0pt}
\setlength{\parskip}{5pt}
\setlength{\emergencystretch}{3em}
\newcommand{\N}{\mathbb{N}}
\newcommand{\Z}{\mathbb{Z}}
\newcommand{\Q}{\mathbb{Q}}
\newcommand{\R}{\mathbb{R}}
\newcommand{\C}{\mathbb{C}}
\newcommand{\F}{\mathbb{F}}
\newcommand{\A}{\mathbb{A}}
\newcommand{\PP}{\mathbb{P}}
\newcommand{\E}{\mathbb{E}}
\newcommand{\eps}{\varepsilon}
\newcommand{\dd}{\,\mathrm{d}}
\begin{document}
\section*{UnsolvedMath ID: 30005821}
\textbf{Source problem code:} catalogue entry
\subsection{Definitions and mathematical statement}
The supplied record combines two questions. They are retained separately, with unresolved conventions stated explicitly.

\begin{enumerate}
\item \textbf{Dihedral-angle gap.} For a fixed dimension \(d\), ask for the largest angle gap \(\varepsilon_d\) such that there exist infinitely many \(d\)-polytopes whose every dihedral angle is at most \(\pi-\varepsilon_d\).

Under the usual interpretation of a \(d\)-polytope as a full-dimensional convex polytope in \(\R^d\), with \(d\geq 2\), a ridge is a face of dimension \(d-2\). Write \(\theta_P(F)\) for the interior dihedral angle, measured in radians, between the two facets incident to a ridge \(F\) of \(P\). The required angle condition is
\[
\theta_P(F)\leq \pi-\varepsilon_d
\qquad\text{for every ridge }F\text{ of every polytope }P\text{ in the family}.
\]
More explicitly, once the intended class of polytopes and the intended notion of distinctness have been specified, define
\[
A_d=\left\{\varepsilon\in\R:\begin{array}{l}
\text{there exists an infinite family }\mathcal P\text{ of pairwise distinct}\
\text{polytopes in that class such that }\theta_P(F)\leq\pi-\varepsilon\\
\text{for every }P\in\mathcal P\text{ and every ridge }F\text{ of }P
\end{array}\right\}.
\]
The question requests \(\max A_d\), provided this maximum exists. A supremum and an attained maximum must be distinguished.

\emph{Unresolved qualifications.} The supplied statement does not specify the intended range of \(d\), convexity, the domain of the gap parameter, or what makes two polytopes distinct. The convex interpretation above is therefore conditional. Distinct positioned subsets, distinct congruence classes, distinct similarity classes, and distinct combinatorial types give different meanings to an infinite family. In particular, translations or dilations preserve all dihedral angles and can yield infinitely many distinct subsets from a single polytope. The intended equivalence convention is essential to the extremal question and cannot be inferred from the supplied record.

\item \textbf{Independent sets in box intersection graphs.} An axis-parallel box in \(\R^3\) has the form
\[
B=I_1\times I_2\times I_3,
\]
where each \(I_j\) is an interval in \(\R\). For a family \(B_1,\ldots,B_n\), its intersection graph \(G\) has vertex set \(\{1,\ldots,n\}\), with an edge between distinct vertices \(i,j\) exactly when \(B_i\cap B_j\neq\varnothing\). Its clique number \(\omega(G)\) is the largest size of a set of pairwise adjacent vertices. Its independence number \(\alpha(G)\) is the largest size of a set of pairwise nonadjacent vertices.

The stated question is whether
\[
\exists c>0\ \forall n\in\N,\ n\geq1\ \forall(B_1,\ldots,B_n),
\qquad
\omega(G)=2\ \Longrightarrow\ \alpha(G)\geq cn.
\]
Here \(c\) is independent of \(n\) and of the box family. Equivalently, the conclusion asks for a set \(S\subseteq\{1,\ldots,n\}\) such that
\[
|S|\geq cn,
\qquad B_i\cap B_j=\varnothing\quad\text{for all distinct }i,j\in S.
\]

\emph{Unresolved qualifications.} The original text glosses clique number \(2\) as ``no three meeting at a common point.'' For axis-parallel boxes, the latter condition corresponds to \(\omega(G)\leq2\), whereas \(\omega(G)=2\) additionally requires an edge. The displayed formulation preserves the equality in the supplied clean statement; the background describes the triangle-free class. The source does not specify whether boxes are closed, open, bounded, or required to have positive side lengths, or how boundary contacts are treated beyond the chosen set-intersection convention.
\end{enumerate}

\textbf{Extraneous source fragments.} The raw original statement continues with incomplete text about a guarantee of at least \(m(n)\) monochromatic empty triangles, the question \(f(n)=o(n^2)\), a claimed lower bound \(f(n)\geq cn^{4/3}\) for a positive constant \(c\), and constructions attributed to B\'ar\'any and Valtr of uncolored \(n\)-point sets having only \(O(n^2)\) empty triangles. It defines an empty triangle as a triangle determined by three of the points whose interior contains no other point. The quantities \(f(n)\) and \(m(n)\), the full point-set hypotheses, and their relationship are absent. The lower-bound claim is attributed by an unidentified speaker to joint work with G\'eza T\'oth. A further fragment, labeled Problem 20 and attributed to Arnau Padrol, asks for the maximum number of vertices of the convex hull of a subset of the \(k\)-barycenters of \(n\) points in \(\R^d\), without defining \(k\)-barycenters or specifying the parameter ranges. These fragments are preserved as source limitations; they do not supply sufficiently complete additional targets here.
\subsection{Short English statement}
For permanent UnsolvedMath ID \(30005821\), the supplied record asks two questions: determine the largest uniform gap below \(\pi\) for all dihedral angles in an infinite family of \(d\)-polytopes, and decide whether every \(n\)-vertex box intersection graph in \(\R^3\) with clique number \(2\) has an independent set of size at least \(cn\) for one universal constant \(c>0\). The first question requires clarification of how polytopes are counted. Box conventions and incomplete unrelated source fragments also require review. The supplied open label is dated curation, not an independently verified status.
\subsection{Sources}
\begin{itemize}
\item Local registry: \texttt{data/unsolvedmath/problems.json}, permanent UnsolvedMath ID \(30005821\), title ``Dihedral-Angle Gaps and Box-Graph Independence,'' exact source code \texttt{catalogue entry}, Oberwolfach Reports Open Problems collection, proposed year \(2024\). No proposer is supplied for the two principal questions.
\item Supplied primary-report citation: \emph{Discrete Geometry} (2024), Oberwolfach Reports, DOI \(10.4171/owr/2024/3\), pp.~137--202. \url{https://doi.org/10.4171/owr/2024/3}. The supplied background describes these reports as records of mathematical workshops and their open-problem discussions.
\item Supplied literature source: ``Oberwolfach Report 3/2024.'' \url{https://ems.press/content/serial-article-files/48651}.
\item Supplied literature source: Walczak, \emph{Triangle-free geometric intersection graphs with no large independent sets}. \url{https://doi.org/10.1007/s00454-014-9645-y}. The record does not state a precise theorem from this reference or explain its exact implications for the second question.
\item Reference supplied within the extraneous empty-triangle fragment: J.~Pach and G.~T\'oth, \emph{Monochromatic empty triangles in two-colored point sets}, Discrete Applied Mathematics \textbf{161}(9), 1259--1261 (2013). No URL is supplied for this reference. The same fragment attributes uncolored empty-triangle constructions to B\'ar\'any and Valtr without a bibliographic citation, and attributes its separate Problem 20 to Arnau Padrol.
\item Literature assessment date: \texttt{2026-08-22}. The supplied dated curation labels the record open and asserts that the report and subsequent literature do not resolve either principal statement. It describes the record as concatenating two report problems and calls the assessment dated literature triage requiring independent verification of time-sensitive status and cited claims. This is supplied dated curation, not independently verified here.
\item The registry assigns statement status \texttt{corrected\_verified} and says that the statement was reconstructed and checked against the official Oberwolfach report. This verification claim is supplied metadata; no source visit or independent verification is asserted in this draft.
\end{itemize}
\end{document}
